% Generated from ../chapters/15-what-practice-means.md
\chapter{Chapter 15 --- What Practice Means}\label{chapter-15-what-practice-means}

Only now can we define the ``practitioner'' addressed by this book.

A practitioner is not necessarily a follower of one school. The practice is the deliberate cultivation of adaptive capacity across levels of the organism. Tai Chi, Qigong, meditation, strength training, walking, difficult intellectual work, relationships, sleep, and nutrition can all contribute, but none is sacred.

The practice has four recurring movements.

First, \textbf{attention}. The person learns to observe without immediately forcing an interpretation or correction. Attention increases the resolution of feedback.

Second, \textbf{release}. Unnecessary muscular, cognitive, and identity constraints are allowed to soften. Release is not collapse. It preserves safety and intention while reducing micromanagement.

Third, \textbf{challenge}. The organism encounters demands slightly beyond its current stable repertoire. Physical training challenges force and coordination. An ambiguous project challenges models, beliefs, and identity. Social responsibility challenges emotional regulation and purpose.

Fourth, \textbf{integration}. Sleep, nourishment, rest, low-intensity movement, and time allow adaptation to become reproducible rather than episodic.

These movements form a cycle rather than a checklist. Challenge without release can deepen compensation. Release without challenge can become passive comfort. Recovery without adequate stimulus produces no new capacity. Stimulus without recovery produces decline.

The guiding principle is not maximal stress but maximum \textbf{integrable} challenge: the largest demand from which the whole organism can recover with greater future capacity.
